% !TeX root = ../main.tex

\setproblem{K}{\ProblemTitleK}
\section{\texorpdfstring{\pletter. \ptitle}{문제 \pletter}}

\begin{frame}[label=problem-K]
    \sectiontitle{\pletter}{\ptitle}
    \sectionmeta{
        출제자: \idtext{ehdnjs14}\\
        태그: \tagtext{\#greedy}\\
        예상 난이도: \tiertext{Silver4}
    }
    \begin{itemize}
        \item 제출 19번, 정답 6명 (정답률 31.6\%)
        \item 처음 푼 사람: ㅅㅇㅅ (00:34:28)
    \end{itemize}
\end{frame}

\begin{frame}{\textbf{\pletter}. \ptitle}
    \textbf{$3$개짜리부터 최대한 고르기}
\begin{itemize}
        \item 얼음이 $1,2,3$개인 정수기의 수를 각각 $c_1,c_2,c_3$에 셉니다. $0$개짜리는 고르지 않습니다.
        \item $K$개를 넘기지 않으면서 고를 수 있는 $3$개짜리의 최대 대수를 $g$라 합시다.
        \[
            g=\min(c_3,\lfloor K/3\rfloor)
        \]
        \item 먼저 $3$개짜리를 $g$대 고르고, 남은 얼음을 $1,2$개짜리로 채워 봅시다.
        \item 다만 $K=4$이고 정수기에 $3,2,2$개가 들어 있다면, $3$개짜리를 빼야 $2+2$로 채울 수 있습니다.
        \item 이처럼 한 대를 빼는 경우까지 확인하면 충분합니다. 최적해는 $g$대 또는 $g-1$대를 사용합니다.
        \item $g=0$이면 $3$개짜리를 사용하지 않는 경우만 확인합니다.
    \end{itemize}
\end{frame}

\begin{frame}{\textbf{\pletter}. \ptitle}
    \textbf{두 경우만 확인해도 되는 이유}
\begin{itemize}
        \item 최적해에서 $3$개짜리를 $x\le g-2$대 사용한다고 가정해 봅시다.
        \item 사용하지 않은 $3$개짜리가 최소 두 대 있고, $1,2$개짜리로 얻는 얼음은 $K-3x\ge6$개입니다.
        \item 이때 사용한 $1,2$개짜리 중 합이 $6$인 정수기들을 항상 고를 수 있습니다.
        \item $2$개짜리가 세 대 이상이면 세 대를 고릅니다. 그보다 적으면 모두 고른 뒤 $1$개짜리로 채웁니다.
        \item 이렇게 고른 정수기는 최소 세 대입니다. 이를 $3$개짜리 두 대로 바꾸면 얼음 수는 같고 사용 대수는 줄어듭니다.
        \item 최적해라는 가정에 모순이므로, $3$개짜리를 $g-2$대 이하로 사용하는 경우는 확인하지 않아도 됩니다.
    \end{itemize}
\end{frame}

\begin{frame}{\textbf{\pletter}. \ptitle}
    \textbf{남은 얼음을 $1,2$개짜리로 채우기}
\begin{itemize}
        \item $3$개짜리의 사용 대수 $x$를 정하면 남은 얼음은 $r=K-3x$개입니다.
        \item $2$개짜리를 최대한 쓰고, 나머지를 $1$개짜리로 채웁니다. 사용 대수는 각각 $y,z$입니다.
        \[
            y=\min(c_2,\lfloor r/2\rfloor),\qquad z=r-2y
        \]
        \item 사용 대수는 $x+y+z=x+r-y$이므로, $y$가 클수록 유리합니다.
        \item $z>c_1$이면 불가능합니다. $2$개짜리를 덜 쓰면 $1$개짜리가 더 많이 필요합니다.
        \item $x=g,g-1$ 중 음수가 아닌 경우만 확인해, 가능한 사용 대수의 최솟값을 구합니다.
        \item 두 경우 모두 불가능하면 $-1$을 출력합니다.
        \item 정수기를 센 뒤 두 경우만 확인하므로 시간 복잡도는 $O(N)$입니다.
    \end{itemize}
\end{frame}
